\chapter{Referência Rápida de Comandos}

As tabelas reúnem a interface pública mais usada. Nos exemplos, argumentos entre colchetes são opcionais; argumentos entre chaves são obrigatórios.

\section{Informações gerais}

\begin{longtable}{>{\ttfamily\footnotesize\raggedright\arraybackslash}p{0.40\textwidth}p{0.20\textwidth}p{0.30\textwidth}}
\toprule
\normalfont\textbf{Comando} & \textbf{Finalidade} & \textbf{Uso e observações} \\
\midrule
\endhead
\string\autor & Autor único & Parâmetro: nome. Use \string\autores para trabalho coletivo. \\
\string\autores & Dois ou mais autores & Parâmetro: nomes separados por \texttt{//}. \\
\string\titulo & Título principal & Parâmetro: texto sem a pontuação do subtítulo. \\
\string\subtitulo & Subtítulo opcional & Parâmetro: texto; deixe vazio quando inexistente. \\
\string\curso & Curso & Parâmetro: nome por extenso. \\
\string\tipotrabalho & Tipo documental & Parâmetro: tipo do trabalho. \\
\string\graupretendido & Grau pretendido & Parâmetro: texto que completa a natureza. \\
\string\local & Local & Parâmetro: cidade. \\
\string\ano & Ano & Parâmetro: quatro algarismos. \\
\string\orientador & Orientador & Parâmetro: nome sem prefixo acadêmico. \\
\string\coorientador & Coorientador & Parâmetro: nome; pode ficar vazio. \\
\string\areadeconcentracao & Área de concentração & Parâmetro: área, quando pertinente. \\
\string\dataaprovacao & Data de aprovação & Parâmetro: local e data por extenso. \\
\string\membrobanca & Membro da banca & Parâmetros: nome, prefixo e instituição. \\
\bottomrule
\end{longtable}

\section{Conteúdo e arquivos}

\begin{longtable}{>{\ttfamily\footnotesize\raggedright\arraybackslash}p{0.40\textwidth}p{0.20\textwidth}p{0.30\textwidth}}
\toprule
\normalfont\textbf{Comando} & \textbf{Finalidade} & \textbf{Uso e observações} \\
\midrule
\endhead
\string\undfsetcontent & Registrar conteúdo & Parâmetros: chave e conteúdo. \\
\string\undfusecontent & Recuperar conteúdo & Parâmetro: chave. \\
\string\undfifcontent & Testar conteúdo & Parâmetros: chave, saída preenchida e saída vazia. \\
\string\undfinputoptional & Incluir arquivo opcional & Parâmetro: caminho sem \texttt{.tex}. \\
\string\undfincludechapter & Incluir capítulo & Parâmetro: caminho sem \texttt{.tex}. \\
\string\undfwarnifempty & Avisar ausência & Parâmetros: chave e mensagem; não interrompe. \\
\string\undfincluirpdf & Incluir PDF integral & Parâmetro: arquivo; aceita opções em colchetes. \\
\bottomrule
\end{longtable}

\section{Elementos pré e pós-textuais}

\begin{longtable}{>{\ttfamily\footnotesize\raggedright\arraybackslash}p{0.40\textwidth}p{0.22\textwidth}p{0.28\textwidth}}
\toprule
\normalfont\textbf{Comando} & \textbf{Resultado} & \textbf{Observações} \\
\midrule
\endhead
\string\undfcover & Capa & Usa as informações gerais. \\
\string\undffolhaderosto & Folha de rosto & Usa natureza e orientação. \\
\string\undfversofolhaderostoopcional & Ficha no verso & Só aparece com conteúdo configurado. \\
\string\undferrata & Errata & Só aparece quando preenchida. \\
\string\undffolhadeaprovacao & Folha de aprovação & Usa data e membros da banca. \\
\string\listadeilustracoes & Lista de figuras & Gerada a partir das legendas. \\
\string\listadefluxogramas & Lista própria & Exibida quando a contagem exige. \\
\string\listadetabelas & Lista de tabelas & Gerada a partir das legendas. \\
\string\undfsumario & Sumário & Exige novas passagens após mudanças. \\
\string\printundfreferences & Referências & Usa a base ligada por \string\addbibresource. \\
\string\undfapendice & Apêndice & Parâmetros: título e rótulo. \\
\string\undfanexo & Anexo & Parâmetros: título e rótulo. \\
\string\undfapendicepdf & Apêndice em PDF & Parâmetros: título e arquivo. \\
\string\undfanexarpdf & Anexo em PDF & Parâmetros: título e arquivo. \\
\bottomrule
\end{longtable}

\section{Citações e referências}

\begin{longtable}{>{\ttfamily\footnotesize\raggedright\arraybackslash}p{0.40\textwidth}p{0.22\textwidth}p{0.28\textwidth}}
\toprule
\normalfont\textbf{Comando} & \textbf{Resultado} & \textbf{Erro comum} \\
\midrule
\endhead
\string\citarparenteses & Citação parentética & Localizador opcional e chave obrigatória. \\
\string\citartexto & Autor integrado ao texto & Localizador opcional e chave obrigatória. \\
\string\citardireta & Citação curta parentética & Localizador, chave e texto. \\
\string\citardiretatexto & Citação curta narrativa & Localizador, chave e texto. \\
\string\citacaolonga & Bloco de citação longa & Localizador opcional e chave obrigatória. \\
\string\parenciteapud & Citação de citação & Autoria/ano originais e chave da consultada. \\
\string\textapud & Autoria integrada ao texto & Mesmos dados e localizadores opcionais. \\
\string\apudbib & Interface compatível anterior & Dois localizadores e duas chaves. \\
\string\fullref & Referência completa pontual & Parâmetro: chave. \\
\string\notarodape & Nota de rodapé & Parâmetro: texto. \\
\bottomrule
\end{longtable}

\section{Objetos e referências cruzadas}

\begin{longtable}{>{\ttfamily\footnotesize\raggedright\arraybackslash}p{0.43\textwidth}p{0.20\textwidth}p{0.27\textwidth}}
\toprule
\normalfont\textbf{Comando ou ambiente} & \textbf{Resultado} & \textbf{Observações} \\
\midrule
\endhead
\string\undfsecao & Seção rotulada & Parâmetros: título e rótulo. \\
\string\undfsubsecao & Subseção rotulada & Parâmetros: título e rótulo. \\
\string\equacao & Equação numerada & Rótulo opcional e expressão. \\
\string\undffigura & Figura & Largura, arquivo, legenda, rótulo e fonte. \\
\string\undffigurasladoalado... & Duas figuras independentes & Cada figura recebe seus próprios dados. \\
\string\undftabela & Tabela flutuante & Legenda, rótulo e fonte opcional. \\
\string\undffluxograma & Fluxograma & Legenda, rótulo e fonte opcional. \\
\string\undfpseudocodigo & Pseudocódigo comum & Assinatura, rótulo e título. \\
\string\undfpseudocodigolongo & Pseudocódigo multipágina & Mesmos parâmetros; permite quebra. \\
\string\figuraref & Chamada de figura & Parâmetro: rótulo sem \texttt{fig:}. \\
\string\tabelaref & Chamada de tabela & Parâmetro: rótulo sem \texttt{tab:}. \\
\string\equacaoref & Chamada de equação & Parâmetro: rótulo sem \texttt{eq:}. \\
\string\pseudocodigoref & Chamada de pseudocódigo & Parâmetro: rótulo sem \texttt{alg:}. \\
\bottomrule
\end{longtable}

\section{Validações}

\begin{longtable}{>{\ttfamily\footnotesize\raggedright\arraybackslash}p{0.40\textwidth}p{0.24\textwidth}p{0.26\textwidth}}
\toprule
\normalfont\textbf{Comando} & \textbf{Efeito} & \textbf{Observação} \\
\midrule
\endhead
\string\undfvalidateonepagecontents & Executa estimativas & Chame antes da capa. \\
\string\bloquearresumoslongos & Excesso vira erro & Opção de revisão, não padrão. \\
\string\avisarresumoslongos & Excesso permanece aviso & Comportamento padrão. \\
\string\undfdesabilitaravisosimagem & Suspende avisos de peso & Não altera a imagem. \\
\string\undfhabilitaravisosimagem & Reativa avisos & Ativo por padrão. \\
\bottomrule
\end{longtable}

\section{Exemplos e resultados esperados}

As duas tabelas seguintes se complementam. A primeira apresenta finalidade, parâmetros e uso básico; a segunda registra resultado esperado, erro comum e solução. Os comandos de apoio menos frequentes permanecem documentados nas tabelas temáticas anteriores.

\begin{landscape}
\begin{longtable}{>{\ttfamily\footnotesize\raggedright\arraybackslash}p{0.25\linewidth}p{0.19\linewidth}p{0.22\linewidth}>{\ttfamily\footnotesize\raggedright\arraybackslash}p{0.26\linewidth}}
\toprule
\normalfont\textbf{Comando} & \textbf{Finalidade} & \textbf{Parâmetros} & \normalfont\textbf{Uso básico} \\
\midrule
\endhead
\string\autor & Registrar autoria & Nome obrigatório & \string\autor\{Ana Lima\} \\
\string\citarparenteses & Inserir citação parentética & Localizador opcional e chave obrigatória & \string\citarparenteses\{silva2025\} \\
\string\undffigura & Inserir figura & Largura opcional, arquivo, título, rótulo e fonte opcional & \string\undffigura\{img\}\{Título\}\{id\} \\
\string\equacao & Inserir equação & Rótulo opcional e expressão obrigatória & \string\equacao[e]\{E=mc\string^2\} \\
\string\undfsecao & Criar seção rotulada & Título e rótulo obrigatórios & \string\undfsecao\{Objetivos\}\{obj\} \\
\string\bloquearresumoslongos & Ativar erro optativo de revisão & Nenhum & \string\bloquearresumoslongos \\
\bottomrule
\end{longtable}

\begin{longtable}{>{\ttfamily\footnotesize\raggedright\arraybackslash}p{0.25\linewidth}p{0.23\linewidth}p{0.21\linewidth}p{0.23\linewidth}}
\toprule
\normalfont\textbf{Comando} & \textbf{Resultado esperado} & \textbf{Erro comum} & \textbf{Solução} \\
\midrule
\endhead
\string\autor & Autoria nas páginas iniciais & Deixar o exemplo padrão & Substituir pelo nome real antes de compilar \\
\string\citarparenteses & Chamada autor-data vinculada à bibliografia & Usar chave inexistente & Conferir a chave no arquivo bibliográfico e executar o Biber \\
\string\undffigura & Figura numerada, com legenda, rótulo e fonte & Informar caminho incorreto & Conferir nome, extensão e caminho relativo do arquivo \\
\string\equacao & Equação numerada e referenciável & Repetir o rótulo & Usar um identificador único e recompilar \\
\string\undfsecao & Seção numerada e referenciável & Repetir o rótulo & Adotar rótulo único e descritivo \\
\string\bloquearresumoslongos & Excesso estimado tratado como erro & Ativar como comportamento comum de escrita & Reservar o comando para revisão controlada; o padrão oficial é aviso \\
\bottomrule
\end{longtable}
\end{landscape}
